\section{Source guide and limits of coverage}\label{sec:sources}
The selection starts from the two requested articles,
\href{https://en.wikipedia.org/wiki/Graph_(discrete_mathematics)}{\emph{Graph
(discrete mathematics)}} and \href{https://en.wikipedia.org/wiki/Graph_theory}{\emph{Graph
theory}}. The left column identifies an article or a topic reached from their
text. Named corollaries and quantitative variants in the right column are
proved here. Their particular algorithms and constants are not attributed to
the encyclopaedia.

\begin{longtable}{@{} >{\raggedright\arraybackslash}p{0.32\linewidth}@{\hspace{0.04\linewidth}}p{0.62\linewidth}@{}}
\toprule
Topic or immediate link & Statements treated here\\ \midrule\endhead
Definition, directed graphs, regular graphs,
\href{https://en.wikipedia.org/wiki/Degree_(graph_theory)}{degree}
& Incidence counts with loops, parity, degree bounds, graph codes, degree-sequence realization and the six-vertex example.\\[4pt]
\href{https://en.wikipedia.org/wiki/Connectivity_(graph_theory)}{Connectivity},
\href{https://en.wikipedia.org/wiki/Directed_acyclic_graph}{directed acyclic graph}
& Reachability and separation, metric, components, topological order, transitive reduction, DAG enumeration, and the finite/infinite boundary.\\[4pt]
\href{https://en.wikipedia.org/wiki/Tree_(graph_theory)}{Tree},
\href{https://en.wikipedia.org/wiki/Path_graph}{path graph},
\href{https://en.wikipedia.org/wiki/Cycle_graph}{cycle graph},
\href{https://en.wikipedia.org/wiki/Matroid}{matroid}
& Forest characterizations and exchange, cycle and cut bases, Kruskal's algorithm, Cayley's formula, the two matrix--tree identities and their forest versions.\\[4pt]
\href{https://en.wikipedia.org/wiki/Seven_Bridges_of_K%C3%B6nigsberg}{Seven Bridges},
route problems, network flow
& Euler criteria with multiplicity and direction, route inspection, last-exit enumeration, integral max-flow min-cut and edge/vertex Menger theorems.\\[4pt]
\href{https://en.wikipedia.org/wiki/Bipartite_graph}{Bipartite graph},
\href{https://en.wikipedia.org/wiki/Perfect_graph}{perfect graph}
& Odd-cycle alternative, Hall, K\H{o}nig, edge colouring, four perfect families, and stable matching. Finite order-decomposition results are deductions from matching.\\[4pt]
\href{https://en.wikipedia.org/wiki/Graph_operations}{Graph operations},
\href{https://en.wikipedia.org/wiki/Line_graph}{line graph},
\href{https://en.wikipedia.org/wiki/Cartesian_product_of_graphs}{Cartesian product}
& Cartesian, tensor, strong and lexicographic product counts, line graph degrees, complement, union, join, contraction and minors.\\[4pt]
Enumeration, finite random graphs,
\href{https://en.wikipedia.org/wiki/Reconstruction_conjecture}{reconstruction}
& Labelled and unlabelled counts, component recurrence, fixed-pair orbits, deck double counts, regular and disconnected reconstruction, finite probabilistic existence arguments.\\[4pt]
\href{https://en.wikipedia.org/wiki/Petersen_graph}{Petersen graph},
\href{https://en.wikipedia.org/wiki/Cayley_graph}{Cayley graph},
algebraic graph theory
& Automorphism groups, finite Frucht realization, exact Petersen eigenbasis and tree count, adjacency powers, Laplacian kernel and the characteristic-polynomial bipartiteness test.\\[4pt]
\href{https://en.wikipedia.org/wiki/Planar_graph}{Planar graph},
\href{https://en.wikipedia.org/wiki/Dual_graph}{dual graph},
graph polynomials
& Euler characteristic from face splices, density bounds, five colours, dual trees, one direction of the forbidden-minor criterion, chromatic and Tutte specializations, planar flow duality and Tait's three-edge-colouring equivalence.\\[4pt]
\href{https://en.wikipedia.org/wiki/Graph_coloring}{Graph colouring},
\href{https://en.wikipedia.org/wiki/Chordal_graph}{chordal graphs},
\href{https://en.wikipedia.org/wiki/Cograph}{cographs}
& Greedy colouring, Brooks and Vizing, chordal elimination, cograph perfection, chromatic coefficients and Stanley's acyclic-orientation evaluation.\\[4pt]
\href{https://en.wikipedia.org/wiki/Extremal_graph_theory}{Extremal graph theory},
Ramsey theory, Hamiltonian routes
& Mantel and Tur\'an with equality, Erd\H{o}s--Stone with a PR modulus, finite Ramsey bounds, equitable regularity, high girth with high chromatic number, Ore, Dirac and tournaments.\\
\bottomrule
\end{longtable}

\paragraph{Reference points.}
Reinhard Diestel's \href{https://diestel-graph-theory.com/}{\emph{Graph Theory},
sixth edition} gives textbook terminology for these topics. Jeff Erickson's
\href{https://jeffe.cs.illinois.edu/teaching/comptop/2023/notes/09-planar-graphs.html}{notes on planar graphs}
describe the dart convention and its interpretation as an embedding.
The arithmetic presentation and the proof interfaces come from
\repo{README.md}{the ProveIt overview}, its
\repo{ConciseIPRA/Lean/ConciseIPRA/PR.lean}{ordinary iPRA syntax and axioms},
and its \repo{Derivations/Lean/IPRADerivations/Core/Fubini.lean}{double-sum}
and \repo{Derivations/Lean/IPRADerivations/Core/SearchParameters.lean}{bounded-search}
sources. Repository links are pinned to commit \texttt{bbb85238dac5}.
The supplied paper is self-contained for its graph-theoretic arguments.

\paragraph{Results beyond this collection.}
The five-colour construction does not establish the four-colour theorem.
The minor argument proves that a spherical graph excludes $K_5$ and $K_{3,3}$
minors. Proving that every graph excluding them has a spherical rotation
would require the converse direction of Wagner's theorem. The spectral
statements use exact polynomial coefficients and supplied rational vectors.
A general analytic spectral theorem would need a representation for
algebraic or real eigenvalues. The reconstruction proof covers the two
classes stated in Theorem~\ref{thm:deck}. No general reconstruction assertion
is used. These are boundaries of the proofs in this paper, not
unprovability assertions about iPRA.
